%TeX-definitions
\usepackage{hyperref} 
\usepackage{amsmath, array, amssymb}
\usepackage{amsthm}
\usepackage{mathdots,caption}

\usepackage{paralist}
\usepackage[mathcal]{eucal}
\usepackage{booktabs,array,tabularx}
\newcommand{\otoprule}{\midrule[\heavyrulewidth]}
\renewcommand{\tabularxcolumn}[1]{>{\arraybackslash}m{#1}}
\usepackage{framed}
\usepackage[x11names]{xcolor}
\definecolor{shadecolor}{gray}{0.9}
\usepackage[british]{babel}

\usepackage{enumerate}
\usepackage{version}


\renewcommand{\labelenumi}{\bf\arabic{enumi}.}
\renewcommand{\labelenumii}{\bf\alph{enumii}}
\usepackage{multicol,mdwlist}

\newtheoremstyle{oefening}
{\baselineskip}%	<Space above> 
{}%	<Space below> 
{}          %	<Body font>
{}                %	<Indent amount> 
{\bfseries}       %     <Theorem head font> 
{}               %	<Punctuation after theorem head> 
{0.01em}           %	<Space after theorem head>2 
{\thmname{#1}\thmnumber{ #2}} %    <Theorem head spec (can be left empty, meaning `normal')>

\newtheoremstyle{oefening1}
{\baselineskip} % <Space above>
{}              % <Space below>
{}              % <Body font>
{}              % <Indent amount>
{\bfseries}     % <Theorem head font>
{}             % <Punctuation after theorem head>
{}         % <Space after theorem head>
{\thmname{#1}\thmnumber{ #2}\phantom{asdf}\\\protect\vspace{1em}}

\theoremstyle{oefening}
\newtheorem{exercise}{Exercise}

\makeatletter
\@addtoreset{exercise}{section}
\makeatother
% ---------------------------------

\usepackage{subcaption}
%\usepackage{subfigure}

\usepackage[
   width=\textwidth, 
   parskip=5pt, 
   font={small,it}, 
   labelfont={up,bf}
]{caption} 

\usepackage[xetex]{graphicx}


%DOCUMENT LAYOUT
\usepackage{geometry} 
\geometry{
  a4paper,
  left   = 1.25in,
  right  = 1.25in,
  top    = 1in,
  bottom = 2in
}

\usepackage{setspace}
\setstretch{1.1}

\usepackage{parskip}

\raggedbottom

%\setlength\parindent{3em}
%\setlength{\parindent}{1em}


% FONTS
\usepackage{fontspec,xltxtra,xunicode} 
\fontspec[Numbers=OldStyle]{Junicode}

\setromanfont[Mapping=tex-text, Ligatures=Common,
      Numbers=OldStyle]{Junicode}
\setsansfont [Mapping=tex-text]{Arial}
 
\setcounter{secnumdepth}{3}

\usepackage[pagestyles]{titlesec}

\titleformat*{\section}{\huge\bfseries}
\titleformat*{\subsection}{\bfseries}
\titleformat*{\subsubsection}{\itshape}
\titleformat{\paragraph}[runin]{\bfseries}%
   {}{0.25em}{}[.]

\titlespacing*{\section   }   {0pt}{    \baselineskip}{2\baselineskip}
\titlespacing*{\subsection}   {0pt}{0.71\baselineskip}{0.25\baselineskip}
\titlespacing*{\subsubsection}{0pt}{0.5 \baselineskip}{0\baselineskip}
\titlespacing*{\paragraph }   {0pt}{0.36 \baselineskip}{0.25em}

\renewpagestyle{plain}{
  %\global\topskip\normaltopskip
  %\sethead[][\normalsize\scshape\sectiontitle][]
  %        {}{\normalsize\scshape\sectiontitle}{}
  \setfoot [][\thepage][]
          {}{\thepage}{}
}

\pagestyle{plain}


\usepackage{graphicx}
%\pagestyle{myheadings}
\usepackage{rotating}
\usepackage{lscape}

\newtheoremstyle{stelling}
{0.5\baselineskip}%	<Space above> 
{0.5\baselineskip}%	<Space below> 
{\itshape}% \addtolength{\leftskip}{2em}}        %	<Body font>
{}                %	<Indent amount> 
{\bfseries}       %     <Theorem head font> 
{.}               %	<Punctuation after theorem head> 
{.5em}            %	<Space after theorem head>2 
{} %    <Theorem head spec (can be left empty, meaning `normal')>

\newtheoremstyle{beroemdestelling}
{0.5\baselineskip}%	<Space above> 
{0.5\baselineskip}%	<Space below> 
{\itshape}        %	<Body font>
{}                %	<Indent amount> 
{\bfseries}       %     <Theorem head font> 
{.}               %	<Punctuation after theorem head> 
{.5em}            %	<Space after theorem head>2 
{\thmnote{ #3}\thmnumber{ #2}} 
                  % <Theorem head spec 
                  % (can be left empty, meaning `normal')>

\newtheoremstyle{definitie}
{0.5\baselineskip}
{0.5\baselineskip}
{}
{}
{\bfseries}{.}{.5em}{}

\theoremstyle{stelling}
\newtheorem{theorem}{Theorem}[section] 
\newtheorem{stelling}{Stelling}[section] 
\newtheorem*{stelling*}{Stelling}
\newtheorem{proposition}[theorem]{Proposition}
\newtheorem{lemma}[theorem]{Lemma}

\theoremstyle{remark}
\newtheorem{remark}[theorem]{Remark}
\newtheorem{corollary}[theorem]{Corollary}

\theoremstyle{definitie}
\newtheorem{definition}[theorem]{Definition} 
\newtheorem{definitie} [theorem]{Definitie } 
\newtheorem{assumption}[theorem]{Assumption}


\theoremstyle{beroemdestelling}
\newtheorem*{famoustheorem}{Theorem}

\setlength{\unitlength}{1cm}

\newcommand\stut{\rule{0pt}{2.6ex}}
\newcommand\onderstut{\rule[-1.2ex]{0pt}{0pt}}
\newcommand\mand{\quad\text{and}\quad}

\newenvironment{bewijs}
{
  \removelastskip \medskip \noindent
  \textbf{Bewijs}\\
}
{
  \qed 
}

\newenvironment{voorbeeld}
{
  \medskip
  \small
  \noindent
  \textbf{Voorbeeld.}\\
}
{
  \medskip
  \normalsize
}


\newenvironment{hardproof}
{
  \removelastskip \medskip \noindent 
  {\bf Proof} 
  \begin{quote}
  \removelastskip\small
}
{
  \normalsize
  \qed
  \end{quote}
}

\newenvironment{answer}{
    \vspace{2\baselineskip}
    \noindent\ignorespaces
      \begin{snugshade}
      \textit{Solution:\\}} {%
      \end{snugshade}
    \par\noindent\ignorespacesafterend
}


\newcommand{\disp}{\displaystyle}